% Generated draft: requires mathematical and source review.
\documentclass[11pt]{article}
\usepackage[margin=25mm]{geometry}
\usepackage{fontspec}
\setmainfont{Latin Modern Roman}
\usepackage{amsmath,amssymb,mathtools,unicode-math}
\setmathfont{Latin Modern Math}
\usepackage{enumitem,needspace,xurl,hyperref}
\setlength{\parindent}{0pt}
\setlength{\parskip}{5pt}
\setlength{\emergencystretch}{3em}
\newcommand{\N}{\mathbb{N}}
\newcommand{\Z}{\mathbb{Z}}
\newcommand{\Q}{\mathbb{Q}}
\newcommand{\R}{\mathbb{R}}
\newcommand{\C}{\mathbb{C}}
\newcommand{\F}{\mathbb{F}}
\newcommand{\A}{\mathbb{A}}
\newcommand{\PP}{\mathbb{P}}
\newcommand{\E}{\mathbb{E}}
\newcommand{\eps}{\varepsilon}
\newcommand{\dd}{\,\mathrm{d}}
\begin{document}
\section*{UnsolvedMath ID: 20002909}
\textbf{Source problem code:} catalogue entry
\subsection{Definitions and mathematical statement}
Let \(\Lambda\) be the symmetric-function ring over \(\Q\), and let \(\widetilde H_\mu(X;q,t)\) be the modified Macdonald functions. Define \(n(\mu)=\sum_{i\ge1}(i-1)\mu_i\), and define \(\nabla\) by
\[
\nabla\widetilde H_\mu=q^{n(\mu')}t^{n(\mu)}\widetilde H_\mu.
\]
The source identifies the bigraded Frobenius series of diagonal harmonics of degree \(n\) with \(\nabla e_n\), and that of the Garsia--Haiman module indexed by \(\mu\vdash n\) with \(\widetilde H_\mu\). For a bigraded \(S_n\)-module \(M\), this series is \(\sum_{a,b}q^at^b\operatorname{ch}(M_{a,b})\), where \(\operatorname{ch}\) is the Frobenius characteristic. The problem asks for meaningful relationships between these modules or series and \(k\)-Schur functions, for positive integers \(k\). It is a research question rather than a single asserted identity or filtration.

For the Garsia--Haiman part, the supplied aid cites Kato's Corollary 10.4: for each \(k\)-bounded partition \(\mu\), meaning \(\mu_1\le k\), a specified pullback of the Garsia--Haiman module has a genuine filtration by Chen--Haiman modules whose characters are modern \(q\)-parameter \(k\)-Schur functions. The precise pullback, module category, grading and character conventions are not supplied here. In particular, this must not be read as an unqualified filtration of the original bigraded \(S_n\)-module, or as an identification of every historical \(k\)-Schur parameter convention. The aid reports refined Macdonald positivity in that specified convention.

For diagonal harmonics the aid reports, for \(n\ge2\),
\[
[h_n]\bigl((\nabla e_n)|_{q=t=1}\bigr)=(-1)^{n-1}.
\]
Here \([h_n]\) means the coefficient of the complete-basis element \(h_{(n)}\) in the expansion in the basis \(\{h_\lambda\}_{\lambda\vdash n}\), not a Schur coefficient and not an unspecified scalar product. If \(k<n\), every element of \(\Lambda_{(k)}=\Q[h_1,\ldots,h_k]\) has zero such coefficient in degree \(n\). Consequently there can be no finite direct identity
\[
\nabla e_n=\sum_{\substack{\lambda\vdash n\\\lambda_1\le k}}a_\lambda(q,t)\,B_\lambda^{(k)}(X;q,t)
\]
when each coefficient is regular at \((1,1)\), each basis function is regular there, and \(B_\lambda^{(k)}(X;1,1)\in\Lambda_{(k)}\), as for the \(k\)-Schur convention in the reported obstruction. Regular means having a well-defined finite specialization at that point. This excludes positive polynomial direct expansions and character-compatible finite filtrations with those characters. It does not exclude singular-coefficient identities, projections, derived constructions or other indirect relationships. For \(k\ge n\), the aid records the stable Schur regime and ordinary Schur positivity as the stable-range relationship, subject to the stated parameter convention. These observations do not resolve every possible meaning of ``nice relationships.
\subsection{Short English statement}
Seek module-theoretic or symmetric-function links between \(k\)-Schur functions, diagonal harmonics and Garsia--Haiman modules. The supplied aid reports a convention-dependent pullback filtration theorem for the latter and an obstruction to regular direct small-\(k\) expansions of the former, while leaving other relationships possible.
\subsection{Sources}
\begin{itemize}
\item AIM workshop \emph{Generalized Kostka polynomials}, Problem (14), \url{https://aimath.org/WWN/kostka/prob2.pdf}.
\item The research summary cites Kato, Corollary 10.4, and names Chen--Haiman modules; it gives no full publication title, date, URL or complete module and pullback definitions.
\item Local UnsolvedMath registry, permanent ID 20002909, exact code \texttt{catalogue entry}; canonical location \texttt{aim-special-functions-notes.json}, notes[13], dated 14 August 2026; statement labeled exact.
\item Research attempt 1: partial result, problem partially solved at the run, rigor label proved, novelty confidence low; registry status partially solved. The filtration attribution, obstruction and stable-range account are supplied machine-generated research aids, not peer reviewed, requiring independent verification before citation.
\end{itemize}
\end{document}
